\section{Introduction to Gravitational Waves Theory}\label{sec:introduction-to-gravitational-waves-theory}

In this chapter we start the discussion surrounding the theory of gravitational waves.
The main goal of this chapter is to gain the necessary terminology of the theory, and talk about the main ideas.
We start by talking about the wave equation, derived from the weak-field approximations.
Afterward, we would discuss some frameworks, mathematical ideas, and case studies that are important in the topic.

Ultimately the discussion in this chapter can be the groundwork for our research status study\ref{ch:feasibility-study}, which is the main body of this document.


\section{The Wave Equation}\label{sec:derivation-of-the-wave-equation}
We start by the weak-field approximation of a general metric.
What we start with is the assumption that our metric $g_{\mu\nu}$ can be considered as a flat minkowski term and a perturbation term $h_{\mu\nu}$.

\begin{equation}
    g_{\mu\nu} = \eta_{\mu\nu} + h_{\mu\nu} + \mathcal{O}(h^2)\label{eq:weak-field-metric}
\end{equation}

The smaller effects are considered negligible.
Another way to get an intuition with this is that we are far from the source of the curvature, in an asymptotic.
And since we considered the background metric to be minkowski, we're in \textit{asymptotically flat} part of the universe.
This consideration (the weak-field) help us with linearization of Einsteins' Field Equations.

\begin{equation}
    G_{\mu\nu} = -\frac{1}{2} \eta_{\alpha\beta}\partial^{\alpha}\partial^{\beta}\bar{h}_{\mu\nu} + \partial^{\alpha}\partial_{(\mu}\bar{h}_{\nu)\alpha} - \frac{1}{2}\eta_{\mu\nu}\partial^{\alpha}\partial^{\beta}\bar{h}_{\alpha\beta}\label{eq:linearized-einstein-field-tensor}
\end{equation}

Where $\bar h_{\mu\nu} = h_{\mu\nu} - \frac12 \eta_{\mu\nu}h^\alpha_\alpha$ The linearized Einstein tensor is invariant with respect to the infinitesimal coordinate transformation of the form:

\begin{equation}
    h_{\mu\nu}\rightarrow h_{\mu'\nu'} = h_{\mu\nu} + 2\partial_{(\mu}\xi_{\nu)}\label{eq:infinitesimal-transformation-of-perturbed-metric}
\end{equation}

In other words, the weak metric is represented by the equivalence classes of metrics linked by the infinitesimal coordinate transformation above.
If one considers $\xi^\alpha$ as the components of a vector, then the transformation can be written in terms of the Lie derivatives, and the infinitesimal coordinate transformation can be interpreted as an infinitesimal diffeomorphism generated by $\xi$.

Notice that the above transformation is analogous to the gauge transformation for the potentials in electromagnetism.

From~\ref{eq:linearized-einstein-field-tensor} one sees that the last two terms contain the divergence of the metric, imposing the Hilbert gauge (or Lorentz gauge) we get:

\begin{equation}
    \square_{\eta} \bar{h}_{\mu\nu} = -\frac{16 \pi G}{c^4}T_{\mu\nu}\label{eq:equations-of-linearized-gr}
\end{equation}

Some keynotes:
\begin{itemize}
    \item It is always possible to reduce to Hilbert gauge by performing an infinitesimal coordinate transformation.
    \item \ref{eq:equations-of-linearized-gr} is a linear equation for the components of $h$.
    At linear order in $h$, the stress energy tensor does not depend on $h$, so in linear GR one can specify the matter source and solve for the metric.
    \item The Bianchi identity in the weak field simplifies, and it is written as the partial derivative of the Einstein tensor since $\partial$ is the connection associated to $\eta$.
    Hence, the weak field equations imply the conservation of the stress-energy tensor on flat background, and that matter does not backreact on the curvature.
\end{itemize}


\subsection{Some Solutions for Weak-field}\label{subsec:some-solutions-for-weak-Field}

Two important and trivial cases for study are the static source and the no-stress source.
A static matter distribution is modeled by a time-independent stress-energy tensor in the form:

\begin{equation}
    T_{\mu\nu} = \rho t_{\mu}t_{\nu}, \ T_{00} = \rho(x^i), \ T_{0i}= T_{ij} = 0\label{eq:static-source}
\end{equation}

where $t^\mu = (\partial_t)^\mu$ is the vector along the time direction of the global inertial coordinates.
With this prescription, the right hand side of~\ref{eq:equations-of-linearized-gr} is time-independent, thus also the gravitational field must be time independent.

\begin{equation}
    \partial_t\bar h_{\mu\nu} = 0\label{eq:time-independent-metric}
\end{equation}

And the linearized Einstein Field Equations reduce to Poisson equations for the components of the metric field:

\begin{equation}
    \begin{cases}
        \nabla^2\bar h_{\mu\nu} = -16\pi\rho & \mu=\nu=0\\
        \nable^2\bar h_{\mu\nu} = 0, \text{otherwise}
    \end{cases}\label{eq:linearized-answers-static-source}
\end{equation}

The Poisson equations have boundary values equal to zero.
Hence, the solution of the $00$ equation is immediately given in terms of the Newton gravitational potential by formally comparing to Newton's $\nabla^2\phi = 4\pi\rho$.
The solution of the other equations (without source terms) is simply zero.

\begin{align}
    h_{\mu\nu} &= -4\phi t_{\mu} t_{\nu} -\frac12 \eta_{\mu\nu}(4\phi)\\
    g_{\mu\nu} &= \eta_{\mu\nu} + h_{\mu\nu} = \eta_{\mu\nu}\left(1-2\phi \right) - 4\phi t_{\mu}t_{\nu}
\end{align}

or

\begin{equation}
    g = -\left(1 + \frac{2\phi}{c^2}\right)d(ct)^2 + \left(
                                                         1-\frac{2\phi}{c^2}
    \right)\delta_{ij}dx^i dx^j \approx -\left(1+\frac{2\phi}{c^2}\right)dt^2 +\delta_{ij}dx^i dx^j\label{eq:static-source-solution}
\end{equation}

\subsection{No-Stresses Source}\label{subsec:no-stresses-source}

A matter distribution with mass-energy density current vector $J^\mu$ and no stresses is modeled by a stress-energy tensor in the form:

\begin{equation}
    T_{\mu\nu} = 2J_{(\mu}t_{\nu)}-2\rho  t_{\mu} t_{\nu}\label{eq:no-stresses-stress-energy-tensor}
\end{equation}

The linearized equations~\ref{eq:equations-of-linearized-gr} read:

\begin{equation}
    \begin{cases}
        \square \bar h_{0\mu} = -16\pi T_{0\mu}, & \mu = 0,\dots,3\\
        \square \bar h_{ij} = 0, & i,j=1,2,3
    \end{cases}\label{eq:linearized-equations-no-stresses}
\end{equation}

If the spatial component of the metric are assumed time-independent, then they are solution of the boundary value problem with the Poisson equation and boundary values reaching zero ar $r\rightarrow\infty$.
This implies they are zero:

\begin{equation}
    \partial_t \bar h_{ij} = 0\implies
    \begin{cases}
        \nabla^2 \bar h_{ij} = 0 \\
        \bar h_{ij}|_{r\rightarrow\infty} =0
    \end{cases} \implies \bar h_{ij}=0\label{eq:linearized-answers-spatial}
\end{equation}

The linearized equations reduce to only those for $\bar h_{0\mu}$, that are formally equivalent to the Maxwell equations in Lorentz gauge for the field.

\begin{equation}
    A_{\mu} \coloneqq -\frac14 \bar h_{0\mu} = -\frac14 \bar h_{\mu\nu}t^{\nu}\label{eq:gauge-no-stresses-solution}
\end{equation}

Once a solution is found, the metric is given by:

\begin{equation}
    \label{eq:metric-no-stresses-source}
    g_{00} = -1 + 2A_0, \ g_{0i}=4 A_i, \ g_{ij} = 1+2A_0 \delta_{ij}
\end{equation}

Reintroducing the factors $c$, the metric reads:

\begin{equation}
    \label{eq:metric-no-stresses-source-complete}
    g = -\left( 1+\frac{2\phi}{c^2} \right)d(ct)^2 + 4A_i d(ct)dx^i +\left( 1-\frac{2\phi}{c^2} \right)\delta_{ij}dx^i dx^j
\end{equation}

\subsection{Propagation of Gravitational Waves}\label{subsec:propagation-of-gravitational-waves}
The linearized EFE in vacuum are homogeneous wave equations for each component of the metric,

\begin{equation}
    \label{eq:homogeneous-wave-equations-for-metric}
    \square_{\eta}\bar h_{\mu\nu} = \eta^{\alpha\beta}\partial_{\alpha}\partial_{\beta}\bar h_{\mu\nu} = 0
\end{equation}

Solutions to this equation can be constructed by superposition of plane waves with (constant) wave vector $k^\mu = (\omega, \vec k)$ and amplitudes $A_{\mu\nu}$:

\begin{align}
    \bar h_{\mu\nu} &= A_{\mu\nu} \exp\left( i k_\rho x^\rho \right)\\
    \partial_{\mu} \bar h_{\alpha\beta} &= A_{\alpha\beta}\partial_{\mu}(\exp\left( ik_\rho x^\rho\right))
\end{align}

Substituting the plane wave ansatz into the equation one soon finds that the wave vector is null.
Indicating that Gravitational Waves propagate at the speed of light.

The Hilbert gauge translates into the 4 equations that imply the waves are transverse to the direction of propagation.

\begin{equation}
    \label{eq:transverse-to-direction-propagation}
    0 = -\partial^{\alpha}\bar h_{\mu\alpha} = ik^{\mu} A_{\mu\nu}\exp\left(
    i k_{\rho}x^{\rho}
    \right) = ik^{\mu}\bar h_{\mu\alpha} \implies k^{\mu} A_{\mu\nu} =0
\end{equation}

\section{The Cause of Gravitational Waves}\label{sec:the-cause-of-gravitational-waves}
